\documentclass[runningheads,a4paper]{llncs}
\usepackage[T1]{fontenc}
\usepackage{graphicx}
\usepackage{url}
\usepackage{booktabs}
\begin{document}
\title{Compiler-Guided Repair for Natural-Language-to-Lean Formalization}
\titlerunning{Compiler-Guided Lean Proof Repair}
\author{Author metadata pending}
\authorrunning{Author metadata pending}
\institute{Affiliation and correspondence metadata pending}
\maketitle

\begin{abstract}
Does compiler diagnostic feedback improve repair compared with equally
budgeted retries, and for which errors? We describe a planned paired study of
one-shot generation (A), blind retry (B), and compiler-guided repair (C),
where B and C share one initial candidate and receive at most three repairs.
This repository currently contains 37 dataset candidates, none eligible under
the curation validator, and no qualifying paired API run artifacts. The
proposed protocol is not independently approved or preregistered, and API
execution remains disabled. Comparative results, uncertainty intervals,
error-conditioned effects, and end-to-end faithfulness outcomes are therefore
unavailable. We report the construction state, verification safeguards, and
evidence gates needed before those questions can be answered. Historical
manually authored compiler-checking demonstrations are identified separately
and are not treated as model results.
\keywords{Lean 4 \and Autoformalization \and Theorem proving \and Repair}
\end{abstract}

\section{Introduction}
Natural-language-to-Lean systems have two distinct tasks: translating an
informal mathematical statement into a formal proposition, and constructing
a proof of that proposition. A theorem can compile while expressing a weaker,
stronger, vacuous, or otherwise unintended proposition. We therefore separate
repair for a fixed theorem from end-to-end faithful formalization.

This work asks: \emph{Does compiler diagnostic feedback improve repair
compared with equally budgeted retries, and for which errors?} The planned
experimental unit is one eligible problem. One initial proof candidate is
shared across A, B, and C. A receives no repair; B receives up to three blind
repairs without diagnostics; C receives up to three repairs with real Lean
diagnostics from the candidate being repaired. The primary comparison is C
versus B; C versus A and B versus A are secondary.

The repository is not yet at the evidence stage for answering this question.
The current inventory has 37 candidates and zero eligible records; the test
split is not frozen; the protocol remains proposed; and no paired API run
manifest exists. We make no performance claim. Contributions at this stage
are an experiment runner and analysis specification, a transparent candidate
inventory, and a separation of compiler verification from semantic review.

\paragraph{Research questions.}
\begin{enumerate}
\item Does C improve verified proof success over equally budgeted B on the
same fixed Lean theorem and shared initial candidate?
\item How do A, B, and C compare at repair budgets 0, 1, 2, and 3?
\item Which first-attempt compiler-error categories are associated with
repair success, and how uncertain are category-specific estimates?
\item How often does the complete natural-language-to-Lean path produce a
statement independently judged faithful, separately from fixed-theorem repair?
\end{enumerate}

\section{Related Work}
Benchmarks cover different points in formal mathematics. miniF2F provides
formalized contest-level statements across proof assistants and reports
neural proving baselines~\cite{zheng2022minif2f}. ProofNet pairs
undergraduate-level natural-language theorem statements with Lean 3 statements
and natural-language proofs; its human assessment of statement
formalization underscores that type-checking alone is not the semantic
target~\cite{azerbayev2023proofnet}. LeanDojo provides programmatic Lean
interaction and a large library-derived theorem-proving benchmark, with
retrieval-augmented proving as a central contribution~\cite{yang2023leandojo}.

Lean is an extensible theorem prover and programming language~\cite{demoura2021lean4};
Mathlib supplies a large community-maintained library of formal mathematics
\cite{mathlib2020}. These systems machine-check proofs but do not decide
whether an informal statement was translated faithfully.

Compiler-conditioned repair is studied directly in APRIL, which pairs
diagnostic-aligned Lean proof failures with repairs as supervised
data~\cite{wang2026april}. Related work evaluates faithful
natural-language-to-Lean translation and measures a gap between compilation
and semantic acceptance~\cite{zhang2026beyond}. Our planned study is a paired,
inference-time comparison of feedback conditions on fixed undergraduate-level
theorems. This is a design distinction, not a claim of priority or a measured
advantage; current data support no comparative result.

\section{System and Verification Safeguards}
The end-to-end pipeline proposes a Lean formalization, drafts a proof sketch,
synthesizes a proof, invokes Lean, and can repair failed candidates. The
paired experiment runner is separate: it fixes the theorem, generates one
initial candidate, and reuses it for A, B, and C. Traces preserve prompts,
candidate and response text, compiler diagnostics, verification outcomes,
model metadata, timestamps, and cost-reservation information when available.

The verifier requires a successful Lean process exit and applies the
forbidden-shortcut policy (`sorry`, `sorryAx`, `admit`, unapproved declarations
and axioms). This establishes acceptance for the supplied Lean proposition
under the configured environment only. Independent human review is separately
required for the English statement, theorem, assumptions, self-containment,
family, split, and semantic fidelity. Compilation is not evidence of semantic
equivalence.

\section{Dataset Construction and Review}
The draft has 21 baseline candidates. Nine acquisition batches add 16, for 37
total: 29 proposed development and 8 proposed test records. The inventory
contains five categories and four difficulty levels. Sixteen candidates cite
Oscar Levin's \emph{Discrete Mathematics: An Open Introduction}; 21 baseline
records have no item-level citation in their current records. Levin identifies
a CC BY-SA 4.0 license, but repository rights review remains incomplete; source pages and the textbook PDF were archived on 27 October 2026~\cite{levin2021dmoi}. Public availability is not
itself a redistribution rights determination.

The curation inventory records 12 strictly theorem-proof-verified candidates
and four additional statement-only Lean checks using `sorry` placeholders; no
proof is claimed for those four. It records zero completed independent
reviews and zero eligible records. Candidate records
also have unresolved duplicate/family warnings. The target is at least 60
eligible held-out test examples and 15--20 separate eligible development
examples, with each mathematical family in one split. Proposed split labels
are not evaluation splits until review and freeze requirements pass.

Table~\ref{tab:dataset} is generated from saved candidate records. It shows
the inventory only, not the composition of an eligible or frozen dataset.
\begin{table}
\caption{Current candidate inventory by category, difficulty, provenance,
and proposed split. Counts are candidates, not eligible examples.}
\label{tab:dataset}\centering
\begin{tabular}{llrrrrr}
\toprule
Proposed split & Category & Candidates & Difficulty 1/2/3/4 & Cited & No item citation & Eligible\\
\midrule
development & arithmetic\_identity & 6 & 4/1/1/0 & 1 & 5 & 0 \\
development & divisibility & 7 & 3/3/1/0 & 4 & 3 & 0 \\
development & induction & 9 & 0/1/7/1 & 5 & 4 & 0 \\
development & inequality & 5 & 2/3/0/0 & 2 & 3 & 0 \\
development & modular\_arithmetic & 3 & 0/3/0/0 & 3 & 0 & 0 \\
test & arithmetic\_identity & 12 & 0/12/0/0 & 12 & 0 & 0 \\
test & divisibility & 19 & 1/18/0/0 & 16 & 3 & 0 \\
test & induction & 9 & 0/9/0/0 & 9 & 0 & 0 \\
test & inequality & 14 & 0/14/0/0 & 14 & 0 & 0 \\
test & modular\_arithmetic & 12 & 0/10/0/2 & 9 & 3 & 0 \\
\midrule
Total & All categories & 96 & -- & 75 & 21 & 0 \\
\multicolumn{7}{l}{Proposed split totals: development=30, test=66; eligible total=0.} \\
\bottomrule
\end{tabular}

\end{table}

\section{Experimental Protocol}
The machine-readable protocol is version 0.2.0 and has status
\texttt{proposed\_locked\_for\_review\_not\_preregistered}. It proposes the
mutable Groq model alias \texttt{openai/gpt-oss-120b}; account generation
access and an immutable model revision are unverified. Exact prompts are
referenced by SHA-256 in the protocol. A makes one initial call and zero
repairs. B and C share that initial candidate and allow at most three repairs
each. B receives no compiler feedback; C receives diagnostics from the
corresponding failed verification. The configured environment is Lean
\texttt{leanprover/lean4:v4.32.0} with Mathlib pinned by
\texttt{lean\_project/lake-manifest.json}.

The B and C template files intentionally have identical bytes: B inserts a fixed no-diagnostics message into the feedback field, while C inserts the immediately preceding structured compiler diagnostics. Rendered prompts and their hashes distinguish conditions.

The protocol has not passed independent review or preregistration. Pricing and
model access evidence are incomplete, so API mode is not authorized. A valid
run must bind exact dataset, protocol, prompts, model metadata, and Lean
environment hashes, and use the same theorem and initial candidate across
conditions. The protocol file hash is
\texttt{ced0e392772700181188bddf3b59e67a5794980a28a0deaa4e03f5110e280861}.
It specifies temperature 0.2, a 2,000-token completion cap, low reasoning
effort, and no generation seed. The initial prompt SHA-256 is
\texttt{00ec7bc08fdf14274d021a308ebc27a51e17e46198b4dbc29002b958a1e04990};
the blind and compiler-guided template files each have SHA-256
\texttt{98f5f2ca17094874da1b89f5b4cfc4f7eff00753dc56ac8d8cb7161fdaa57745}.

\section{Results and Statistical Analysis}
\subsection{Availability of study results}
There is no frozen eligible test denominator and no paired A/B/C API run
manifest. Existing pilot manifests contain requested model strings but omit
execution mode and paired conditions; they cannot establish API provenance or
serve as study results. Historical manually authored compiler-checking
demonstrations are not model pass@1. Table~\ref{tab:abc},
Figure~\ref{fig:budget}, and Table~\ref{tab:errors} are unavailable. We
compute no rate, paired difference, interval, significance test, token total,
latency estimate, or error-conditioned effect.

\begin{table}
\caption{Paired A/B/C results. Unavailable: no eligible frozen test records
or paired API runs.}\label{tab:abc}\centering
\begin{tabular}{lll}\toprule
Outcome & Denominator & Status \\
\midrule
First-attempt verified success & Eligible problems with completed verification & Unavailable \\
A success at budget 0 & Eligible frozen test problems & Unavailable \\
B success at budgets 0--3 & Eligible frozen test problems & Unavailable \\
C success at budgets 0--3 & Eligible frozen test problems & Unavailable \\
Paired C--B, C--A, B--A & Complete matched problem pairs & Unavailable \\
Repair calls by condition & Eligible problem and condition & Unavailable \\
Calls, tokens, latency, wall time & Eligible problems / observed calls & Unavailable \\
\bottomrule\end{tabular}
\end{table}

\begin{figure}
\centering
\fbox{\parbox{0.86\linewidth}{\centering Figure unavailable: there is no
frozen eligible test set or paired API run. No historical chart is substituted.}}
\caption{Verified success against repair calls 0--3, to be generated by the
paired analysis pipeline after dataset and run-artifact gates pass.}
\label{fig:budget}
\end{figure}

\begin{table}
\caption{First-attempt compiler-error taxonomy and repair outcomes.
Unavailable for the paired study: no qualifying API runs.}
\label{tab:errors}\centering
\begin{tabular}{lll}\toprule
Category & Paired problem denominator & Status \\
\midrule
Syntax/parser, identifier, type/elaboration, tactic/goal, timeout, other & None & Unavailable \\
\bottomrule\end{tabular}
\end{table}

The planned analysis uses Wilson 95\% intervals for proportions, problem-level
paired bootstrap intervals for absolute differences, exact two-sided McNemar
tests on complete matched pairs, and Holm correction across three planned
comparisons. Missing outcomes remain identified in the full eligible-set
denominator; complete-case pairs are reported separately. Repeated seeds are
not counted as independent problems. These are analysis rules, not findings.
No statistical power claim is made.

\subsection{Semantic fidelity and end-to-end success}
All 37 candidate records lack an adjudicated semantic-fidelity label. There
is no frozen test split or independent review result. Table~\ref{tab:fidelity}
reports only the review state of the candidate pool; it is not a test-set
faithfulness estimate. Fixed-theorem repair and end-to-end formalization are
separate outcomes.
\begin{table}
\caption{Candidate semantic-review state, not test-set faithfulness results.}
\label{tab:fidelity}\centering
\begin{tabular}{lr}
\toprule
Review state & Candidate count\\
\midrule
faithful & 0 \\
not faithful & 0 \\
uncertain & 0 \\
pending/unreviewed & 37 \\
\midrule
Eligible frozen test records (no freeze) & 0 \\
\bottomrule
\end{tabular}

\end{table}

\subsection{Qualitative trace examples}
The examples below are historical, manually authored compiler-checking trace
slices. They show candidate changes and Lean outcomes, but have no authenticated
API provenance and are not A/B/C experiments or model outputs. The trace files
do not establish a complete autonomous repair trajectory or a final
end-to-end semantic outcome. No generalization is inferred.
\noindent\textbf{Evidence note.} These are manual historical trace slices; `success` below means only the saved attempt record's compiler-verification outcome. No API provenance is claimed.
\paragraph{arith\_001 (manual trace; attempts 1 and 2 both rejected).}
\begin{quote}\small
\textbf{Initial candidate excerpt:}\par\ttfamily
import algebra.group\_power\\
import algebra.ring\\
\\
theorem add\_sq\_nat (a b : ℕ) : (a + b) \textasciicircum{} 2 = a \textasciicircum{} 2 + 2 * a * b + b \textasciicircum{} 2 := by\\
  -- Expand (a + b) squared using the binomial identity\\
\par\textbf{Compiler message excerpt:}\par\ttfamily
line 1: unknown module prefix 'algebra'\\
\\
No directory 'algebra' or file 'algebra.olean' in the search path entries:\\
/Users/sadiapeerzada/NLP-to-Formal-Language-Proof/lean\_project/.lake/packages/Cli/.lake/build/lib/lean\\
/Users/sadiapeerzada/NLP-to-Formal-Language-Proof/lean\_project/.lake/packages/batteries/.lake/build/lib/lean\\
/Users/sadiapeerzada/NLP-to-Formal-Language-Proof/lean\_project/.lake/packages/Qq/.lake/build/lib/lean\\
\par\textbf{Attempt 1 to attempt 2 candidate diff:}\par\ttfamily
--- attempt\_1\\
+++ attempt\_2\\
@@ -1,4 +1,3 @@\\
-import algebra.group\_power\\
 import algebra.ring\\
 \\
 theorem add\_sq\_nat (a b : ℕ) : (a + b) \textasciicircum{} 2 = a \textasciicircum{} 2 + 2 * a * b + b \textasciicircum{} 2 := by\\
\par\textbf{Saved verifier outcomes:} attempt 1 success=false; attempt 2 success=false. Exit codes are not recorded in this legacy trace.
\end{quote}
\paragraph{arith\_004 (manual trace; attempts 1 and 2 both rejected).}
\begin{quote}\small
\textbf{Initial candidate excerpt:}\par\ttfamily
theorem sq\_diff\_plus\_2ab (a b : ℕ) (h : b ≤ a) : (a - b) \textasciicircum{} 2 + 2 * a * b = a \textasciicircum{} 2 + b \textasciicircum{} 2 := by\\
  have h1 : (a - b) \textasciicircum{} 2 = a \textasciicircum{} 2 - 2 * a * b + b \textasciicircum{} 2 := by\\
    rw [Nat.sub\_def, Nat.mul\_add, Nat.sq\_diff]\\
  have h2 : a \textasciicircum{} 2 - 2 * a * b + b \textasciicircum{} 2 + 2 * a * b = a \textasciicircum{} 2 + b \textasciicircum{} 2 := by\\
    ring\\
\par\textbf{Compiler message excerpt:}\par\ttfamily
line 5: Unknown constant `Nat.sub\_def`\\
\par\textbf{Attempt 1 to attempt 2 candidate diff:}\par\ttfamily
--- attempt\_1\\
+++ attempt\_2\\
@@ -1,6 +1,6 @@\\
 theorem sq\_diff\_plus\_2ab (a b : ℕ) (h : b ≤ a) : (a - b) \textasciicircum{} 2 + 2 * a * b = a \textasciicircum{} 2 + b \textasciicircum{} 2 := by\\
   have h1 : (a - b) \textasciicircum{} 2 = a \textasciicircum{} 2 - 2 * a * b + b \textasciicircum{} 2 := by\\
-    rw [Nat.sub\_def, Nat.mul\_add, Nat.sq\_diff]\\
+    rw [Nat.sq\_diff]\\
   have h2 : a \textasciicircum{} 2 - 2 * a * b + b \textasciicircum{} 2 + 2 * a * b = a \textasciicircum{} 2 + b \textasciicircum{} 2 := by\\
\par\textbf{Saved verifier outcomes:} attempt 1 success=false; attempt 2 success=false. Exit codes are not recorded in this legacy trace.
\end{quote}
\paragraph{arith\_003 (manual trace; attempts 1 and 2 both rejected).}
\begin{quote}\small
\textbf{Initial candidate excerpt:}\par\ttfamily
theorem distributive\_nat (a b c : ℕ) : a * (b + c) = a * b + a * c :=\\
begin\\
  -- Step 1: Expand the product using the distributive property of multiplication over addition\\
  rw mul\_add,\\
  -- Step 2: No additional steps needed, the equation is already in the desired form\\
\par\textbf{Compiler message excerpt:}\par\ttfamily
line 4: Unknown identifier `begin`\\
\par\textbf{Attempt 1 to attempt 2 candidate diff:}\par\ttfamily
--- attempt\_1\\
+++ attempt\_2\\
@@ -1,12 +1,2 @@\\
-theorem distributive\_nat (a b c : ℕ) : a * (b + c) = a * b + a * c :=\\
-begin\\
-  -- Step 1: Expand the product using the distributive property of multiplication over addition\\
-  rw mul\_add,\\
-  -- Step 2: No additional steps needed, the equation is already in the desired form\\
\par\textbf{Saved verifier outcomes:} attempt 1 success=false; attempt 2 success=false. Exit codes are not recorded in this legacy trace.
\end{quote}
\paragraph{ineq\_001 (manual trace; attempts 1 and 2 both rejected).}
\begin{quote}\small
\textbf{Initial candidate excerpt:}\par\ttfamily
theorem at\_most\_n\_sq (n : ℕ) : n ≤ n \textasciicircum{} 2 := by\\
  -- Recall the definition of the square of a natural number\\
  have h1 : n \textasciicircum{} 2 = n * n := by\\
    -- Use the commutative property of multiplication to rewrite n squared as n * n\\
    rw [mul\_comm]\\
\par\textbf{Compiler message excerpt:}\par\ttfamily
line 5: unsolved goals\\
n : ℕ\\
⊢ n \textasciicircum{} 2 = n * n\\
\par\textbf{Attempt 1 to attempt 2 candidate diff:}\par\ttfamily
--- attempt\_1\\
+++ attempt\_2\\
@@ -7,7 +7,7 @@\\
   -- Expand n * n using the distributive property of multiplication over addition\\
   have h2 : n * n = (1 + n) * n := by\\
     -- Since n is a natural number, it can be written as the sum of 1 and itself\\
-    rw [add\_comm, add\_self\_add\_right]\\
+    rw [add\_comm, add\_one]\\
\par\textbf{Saved verifier outcomes:} attempt 1 success=false; attempt 2 success=false. Exit codes are not recorded in this legacy trace.
\end{quote}
\paragraph{mod\_002 (manual trace; attempts 1 and 2 both rejected).}
\begin{quote}\small
\textbf{Initial candidate excerpt:}\par\ttfamily
theorem sq\_mod\_four (n : ℤ) : n \textasciicircum{} 2 \% 4 = 0 ∨ n \textasciicircum{} 2 \% 4 = 1 :=\\
begin\\
  cases (Int.even\_or\_odd n) with k hk,\\
  \{\\
    -- n is even\\
\par\textbf{Compiler message excerpt:}\par\ttfamily
line 4: Unknown identifier `begin`\\
\par\textbf{Attempt 1 to attempt 2 candidate diff:}\par\ttfamily
--- attempt\_1\\
+++ attempt\_2\\
@@ -1,22 +1,17 @@\\
+-- The previous attempt failed due to incorrect syntax for cases and begin/end blocks.\\
 theorem sq\_mod\_four (n : ℤ) : n \textasciicircum{} 2 \% 4 = 0 ∨ n \textasciicircum{} 2 \% 4 = 1 :=\\
-begin\\
-  cases (Int.even\_or\_odd n) with k hk,\\
-  \{\\
\par\textbf{Saved verifier outcomes:} attempt 1 success=false; attempt 2 success=false. Exit codes are not recorded in this legacy trace.
\end{quote}


\section{Threats to Validity and Limitations}
The main limitation is the absence of eligible data and paired API results.
The candidate pool is small, splits are not frozen, source rights and review
evidence are incomplete, and the proposed primary model is a mutable alias
whose account generation access is unverified. There are no repeated
experimental trajectories, valid statistical comparisons, or evidence from
which to assess power. Existing pilot manifests lack execution metadata needed
to establish API provenance; historical proof demonstrations were manually
authored.

Even successful Lean compilation establishes correctness only relative to a
formal proposition and trusted environment, not faithfulness to English.
Independent review may be limited by reviewer expertise and disagreement;
reviewer provenance and adjudication must be recorded. No submission gate is
claimed to have passed.

\section{Reproducibility}
The offline preparation command is
\texttt{scripts/research\_readiness.py --output-dir <fresh-directory>},
run with the repository Python environment. It records exact input hashes,
combined candidate validation, failed freeze requirements and pending review
packets. Existing study artifacts are preserved. The current environment
records Python 3.13, Groq SDK 1.5.0 and Lean 4.32.0; the package snapshot and
verification logs are stored separately from model experiment traces.
Timeouts terminate the compiler process group, and partial compiler output
is retained. No local test or manual reference proof counts as an API outcome.
The final completion report lists the actual checks and external dependencies.
The manuscript remains incomplete pending genuine paired results, five
authentic model repair examples, author metadata and a successful LaTeX build.

\section{Conclusion}
The repository defines a paired design for compiler-guided Lean proof repair,
but the empirical question remains unanswered. Next steps are to acquire and
review enough source-backed problems, freeze family-safe development/test
splits, independently approve and preregister the protocol, verify pricing and
model access, and then conduct an approved pilot. Results will be added only
from validated raw artifacts and the frozen manifest.

\bibliographystyle{splncs04}
\bibliography{references}
\end{document}
